\chapter*{Introduction générale}
\addcontentsline{toc}{chapter}{Introduction générale}
\markboth{Introduction générale}{Introduction générale}

L'industrie du sport a connu une métamorphose radicale au cours de la dernière décennie. Avec l'avènement du Big Data et des capacités de calcul à grande échelle, la prise de décision empirique a laissé place à une approche scientifique et mathématique. Dans le football, sport caractérisé par un haut niveau d'incertitude et une rareté des buts (\textit{faible scoring rate}), anticiper l'issue d'une rencontre est un défi de taille qui fascine autant les chercheurs en intelligence artificielle que les analystes sportifs.

\section*{Contexte et Enjeux}
Traditionnellement, la prédiction sportive reposait sur l'expertise humaine, souvent biaisée par des facteurs émotionnels ou des observations subjectives. Aujourd'hui, grâce aux fournisseurs de données comme API-Football, il est possible de quantifier chaque aspect d'une rencontre : pourcentage de possession, tirs cadrés, précision des passes, etc. L'enjeu majeur est de transformer ce volume massif de données brutes en un signal prédictif exploitable, tout en contournant les pièges statistiques liés au football (notamment la difficulté à prédire un match nul).

\section*{Objectifs du Projet}
L'objectif principal de ce projet de fin d'année est de concevoir, développer et déployer une plateforme complète nommée \textbf{Intelligent Football Prediction Dashboard}. Ce système doit accomplir les tâches suivantes :
\begin{enumerate}
    \item \textbf{Ingestion automatisée} : Collecter et stocker massivement les données historiques et tactiques des ligues de football.
    \item \textbf{Modélisation prédictive} : Développer un modèle de Machine Learning robuste capable de calculer les probabilités de victoire (Home, Draw, Away).
    \item \textbf{Restitution visuelle} : Offrir une interface utilisateur (UI) moderne et ergonomique pour l'analyse décisionnelle.
    \item \textbf{Sécurisation} : Garantir l'intégrité des données et des accès via un système d'authentification par rôles.
\end{enumerate}

\section*{Structure du Rapport}
Afin de relater notre démarche d'ingénierie, ce rapport s'articule autour de quatre chapitres :
\begin{itemize}
    \item Le \textbf{Chapitre 1} présente l'analyse fonctionnelle, le cahier des charges et l'état de l'art.
    \item Le \textbf{Chapitre 2} expose les fondements théoriques des algorithmes de Machine Learning utilisés (XGBoost, SMOTE).
    \item Le \textbf{Chapitre 3} détaille la conception architecturale (UML, Base de données).
    \item Enfin, le \textbf{Chapitre 4} couvre l'intégralité de la mise en œuvre technique : du développement Backend à l'interface Frontend, en passant par la phase exhaustive de tests et de validation.
\end{itemize}
